\section{Objetivos del Proyecto}
\subsection{Objetivo General}
\textbf{Evaluar el efecto de un prototipo de control de acceso con códigos QR
	dinámicos sobre el tiempo de validación, la fila y las incidencias de
	identidad en el acceso peatonal piloto del TESCHA, para aportar evidencia que
	permita decidir si conviene ampliarlo a otros accesos.}

Hoy no existen datos que midan el problema, y una decisión sobre el sistema
debe apoyarse en evidencia. El objetivo pide evaluar el efecto, no demostrar
una mejora: si las mediciones no muestran diferencia, ese también es un
resultado válido.

\subsection{Indicadores}
El efecto se mide con tres indicadores, primero antes del piloto (línea base)
y después durante él, en el horario pico del turno vespertino:
\begin{itemize}
	\item \textbf{Tiempo de validación:} promedio por estudiante, desde que
	      llega al punto de revisión hasta que lo cruza.
	\item \textbf{Fila máxima:} la mayor cantidad de personas esperando en el
	      horario pico.
	\item \textbf{Incidencias de identidad:} credenciales ajenas detectadas.
\end{itemize}
La pregunta que guía el trabajo es: ¿cuál es el efecto de implementar un
sistema automatizado de validación mediante códigos QR dinámicos en la
eficiencia general del control de acceso peatonal de los estudiantes del
TESCHA? Por eficiencia se entiende dejar pasar con rapidez a quienes tienen
derecho a entrar, y solo a ellos.

\subsection{Objetivos Específicos}
Los ocho objetivos específicos se reparten en dos líneas de trabajo
(Tabla~\ref{tab:especificos}).

\begin{table}[tbp]
	\caption{Objetivos Específicos}
	\label{tab:especificos}
	\small\setlength{\tabcolsep}{3pt}
	\begin{tabular}{P{1.2cm}P{7.3cm}P{6.6cm}}
		\toprule
		Clave & Qué se hará & Por qué \\
		\midrule
		\multicolumn{3}{l}{\textit{Desarrollo tecnológico}} \\
		OT1 & Diseñar la arquitectura del sistema: lector, torniquete, servidor, base de datos y panel web. & El acceso actual no deja registro digital y el prototipo necesita una estructura que lo registre. \\
		OT2 & Implementar la generación y validación de códigos QR que cambian cada pocos segundos y están firmados, con pruebas de códigos vencidos, repetidos y alterados. & Un código fijo puede copiarse y compartirse; la firma resiste el reenvío y la modificación del mensaje \parencite{subpratatsavee2015hmac}. \\
		OT3 & Construir el prototipo funcional que integre lector, torniquete, servidor y panel web. & Sin prototipo no hay sistema que evaluar. \\
		OT4 & Verificar el desempeño técnico: tiempo de lectura y de consulta, y lecturas fallidas con distintas luces, brillos de pantalla y distancias. & No se encontró literatura que reporte el tiempo de lectura de un QR en este lector, así que debe medirse con el prototipo. \\
		\multicolumn{3}{l}{\textit{Investigación y análisis de datos}} \\
		OI1 & Analizar la situación actual del acceso con observación de cinco días, encuesta y entrevistas, para fijar la línea base. & No existe ningún registro del flujo de acceso; sin línea base no se puede atribuir ningún cambio al sistema. \\
		OI2 & Estimar la viabilidad y la aceptación entre los estudiantes: celular propio, conexión, batería y disposición a mostrar un QR. & El sistema depende del celular del estudiante y se desconoce si la comunidad lo aceptaría. \\
		OI3 & Interpretar las obligaciones de protección de datos personales aplicables al TESCHA \parencite{edomex2017datos, lgpdppso2025}. & El TESCHA es sujeto obligado de ambas leyes y el prototipo registra la hora de ingreso de cada estudiante. \\
		OI4 & Evaluar el efecto del sistema en los tres indicadores, comparando la línea base con el piloto. & La pregunta pide el efecto, y eso solo se conoce comparando un antes y un después. \\
		\bottomrule
	\end{tabular}
	\tablenote{Elaboración propia con base en el documento de investigación del equipo. OT es desarrollo tecnológico; OI es investigación y análisis de datos.}
\end{table}

Los objetivos de desarrollo producen el sistema que se evalúa. Los de
investigación miden los indicadores y aportan lo necesario para interpretarlos.
Sin los primeros no hay sistema que evaluar; sin los segundos no se sabe si el
sistema sirvió.
